% !TeX encoding = UTF-8
% !TeX spellcheck = pt_BR
% !TeX root = 20231-CODDISC-JOAO-FERREIRA.tex
% ------------------------------------------------------------------------------

% ------------------------------------------------------------------------------
% Template para aulas da UNINASSAU
% Autor: João Ferreira da Silva Júnior
% Versão: 0.4.3 - 2023-02-04
% ------------------------------------------------------------------------------



% ------------------------------------------------------------------------------
% Procurar texto e alterar texto em arquivos, buscar em todo subdiretório
% find . -type f -name "*.tex" -exec sed -i 's/TEXTO ORIGINAL/TEXTO NOVO/g' {} +
% find . -type f -name "*.tex" -exec sed -i 's/Versão: 0.1 - 2020-03-27/Versão: 0.4.3 - 2023-02-04/g' {} +
%
% Procurar padrão no nome e alterar arquivos, buscar em todo subdiretório
% find . -name '*.*' -print0 | xargs -0 -n1 bash -c 'mv "$0" "${0/PADRÃO ANTIGO/SUBSTITUTO}"'
% ------------------------------------------------------------------------------



% ------------------------------------------------------------------------------
\documentclass[
  % ----------------------------------------------------------------------------
  11pt,
  % ----------------------------------------------------------------------------
  %aspectratio=43,
  aspectratio=169,
  % ----------------------------------------------------------------------------
  % Opções do pacote babel
  english,                    % idioma adicional para hifenização
  french,                     % idioma adicional para hifenização
  spanish,                    % idioma adicional para hifenização
  brazil                      % o último idioma é o principal do documento
  ]{beamer}
% ------------------------------------------------------------------------------





% ------------------------------------------------------------------------------
% ------------------------------------------------------------------------------
% ------------------------------------------------------------------------------
% ------------------------------------------------------------------------------
% ------------------------------------------------------------------------------
% Configurações visuais e de tema

% Extend beamer and a0poster for custom sized posters
%\usepackage[%
%  orientation=landscape,
%  size=custom,
%  width=16,
%  height=9,
%  scale=0.45
%  ]{beamerposter}

% Configura o avanço do slide para passos a cada objeto ou parágrafo do frame
\usepackage{beamerthemesplit}

% This command allows you to specify in a quite general way how a covered item should be rendered
\setbeamercovered{transparent}

% Navigation symbols are small icons that are shown on every slide by default
\setbeamertemplate{navigation symbols}{}

% Configura o tema e a cor
\usetheme{AnnArbor}
\usecolortheme{default}
\usefonttheme{default}
% ------------------------------------------------------------------------------
% Definição de cores
\definecolor{black}{rgb}{0.0, 0.0, 0.0}
\definecolor{white}{rgb}{255,255,255}
\definecolor{red}{rgb}{255,0,0}
\definecolor{green}{rgb}{0,255,0}
\definecolor{blue}{rgb}{0,0,255}
% ------------------------------------------------------------------------------
% Define a aparência dos blocks
\setbeamertemplate{blocks}[default]
% ------------------------------------------------------------------------------



% ------------------------------------------------------------------------------
% ------------------------------------------------------------------------------
% ------------------------------------------------------------------------------
% ------------------------------------------------------------------------------
% ------------------------------------------------------------------------------
% Pacotes principais

% accept different input encodings (conversão automática dos acentos)
\usepackage[utf8]{inputenc}

% standard package for selecting font encodings
\usepackage[T1]{fontenc}
% ------------------------------------------------------------------------------



% ------------------------------------------------------------------------------
% ------------------------------------------------------------------------------
% ------------------------------------------------------------------------------
% ------------------------------------------------------------------------------
% ------------------------------------------------------------------------------
% Pacotes extras

% the package provides several macros to adjust boxed content
\usepackage{adjustbox}

% Algorithm2e is an environment for writing algorithms. An algorithm becomes a floating object (like figure, table, etc.).
\usepackage{algorithm2e}

% an extended set of fonts for use in mathematics
\usepackage{amsfonts}

% extension package for LaTeX that provides various features to facilitate writing math formulas and to improve the typographical quality of their output
\usepackage{amsmath}

% math symbols defined by LaTeX package
\usepackage{amssymb}
\usepackage{fontawesome5}

% Create PDF and SVG animations from graphics files and inline graphics
% \usepackage{animate}

% draw dash-lines in array/tabular environments
% \usepackage{arydshln}

% manages culturally-determined typographical (and other) rules, and hyphenation patterns for a wide range of languages
\usepackage[english,portuguese]{babel}
  \newcommand{\ingles}[1]{\foreignlanguage{english}{\emph{#1}}}
  \newcommand{\frances}[1]{\foreignlanguage{french}{\emph{#1}}}
  % \newcommand{\grego}[1]{\foreignlanguage{greek}{\emph{#1}}}
  \newcommand{\grego}[1]{\emph{#1}}
  \newcommand{\espanhol}[1]{\foreignlanguage{spanish}{\emph{#1}}}
  \newcommand{\portugues}[1]{\foreignlanguage{brazil}{\emph{#1}}}

% a symbol (dingbat) font and LaTeX macros for its use
\usepackage{bbding}

% enhances the quality of tables in LATEX, providing extra commands as well as behind-the-scenes optimisation
\usepackage{booktabs}

% adding colors to your text is supported by the color package
\usepackage{color}

% gives you Courier scaled at its normal size as the default Typewriter Font
% \usepackage{courier}

% this package provides user control over the layout of the three basic list environments: enumerate, itemize and description
% http://linorg.usp.br/CTAN/macros/latex/contrib/enumitem/enumitem.pdf
% \usepackage{enumitem}

% enhanced support for graphics
\usepackage{graphicx}

% is used to handle cross-referencing commands in LaTeX to produce hypertext links in the document
\usepackage{hyperref}

% this package gives you easy access to the Lorem Ipsum dummy text
\usepackage{lipsum}

% enables the user to typeset programs (programming code) within LATEX
\usepackage{listings}
  \lstset{postbreak=\raisebox{0ex}[0ex][0ex]{\ensuremath{\color{gray}\hookrightarrow\space}}}
  % Aqui tem um exemplo: https://github.com/abntex/abntex2/wiki/HowToFormatarCodigoFonte
  % https://tex.stackexchange.com/questions/64839/how-to-change-listing-caption

  % Listing -> Algorithm
  \renewcommand{\lstlistingname}{Algoritmo}

  % List of Listings -> List of Algorithms
  \renewcommand{\lstlistlistingname}{Lista de algoritmos}

% the Latin Modern family of fonts
\usepackage{lmodern}

% Multimedia inclusion package with Adobe Reader-9/X compatibility
% http://repositorios.cpai.unb.br/ctan/macros/latex/contrib/media9/doc/media9.pdf
\usepackage{media9}

% subliminal refinements towards typographical perfection
\usepackage{microtype}

% Multimedia inclusion package
% Substituído pelo media9
%\usepackage{movie15}
% \newcommand{\pdfmovie}[4]{\href{run:#1}{\framebox{\parbox[c][#3][c]{#2}{\center #4}}}}

% multicol defines a multicols environment which typesets text in multiple columns (up to a maximum of 10)
\usepackage{multicol}

% create tabular cells spanning multiple rows
\usepackage{multirow}

% provides some new list environments. Itemized and enumerated lists can be typesetted within paragraphs
\usepackage{paralist}

% draws high--quality function plots in normal or logarithmic scaling with a user-friendly interface directly in TeX
\usepackage{pgfplots}

% This package reads tab-separated numerical tables from input and generates code for pretty-printed LATEX-tabulars
\usepackage{pgfplotstable}
  \pgfplotsset{compat=1.14}

% provides the means to draw pie (and variant) charts, using PGF/TikZ
\usepackage{pgf-pie}

% this package simplifies the inclusion of external multi-page PDF documents in LaTeX documents
\usepackage{pdfpages}

% defines new commands \Centering, \RaggedLeft, and \RaggedRight and new environments Center, FlushLeft, and FlushRight, which set ragged text and are easily configurable to allow hyphenation
\usepackage{ragged2e}

% makes the page margins visible and hence
% \usepackage{showframe}

% A comprehensive (SI) units package
\usepackage{siunitx}

% provides hyphenatable spacing out (letterspacing), underlining, striking out, etc., using the TeX hyphenation algorithm to find the proper hyphens automatically
  \usepackage{soul}
  \makeatletter
  \let\HL\hl
  \renewcommand\hl{%
    \let\set@color\beamerorig@set@color
    \let\reset@color\beamerorig@reset@color
    \HL}
  \makeatother

% permit use of UTF-8 characters in soul
\usepackage{soulutf8}

% inclusão de Subfiguras https://www.ctan.org/pkg/subfig
\usepackage{subfig}

% for creating graphics programmatically
\usepackage{tikz}
  \usetikzlibrary{arrows,automata,shadows}

% select Adobe Times Roman (or equivalent) as default font
% \usepackage{times}

% Using a color stack for transparency with pdfTEX
\usepackage{transparent}

% provides many glyphs like male and female symbols and astronomical symbols, as well as the complete lasy font set and other odds and ends
\usepackage{wasysym}

% font handwrite
\usepackage{wedn}
  % inline font modify
  \newcommand{\setfont}[2]{{\fontfamily{#1}\selectfont #2}}
  % Usage
  %\setfont{wedn}{Texto...}

% adding colors to your text is supported by the color package
\usepackage{xcolor}
% ------------------------------------------------------------------------------



% ------------------------------------------------------------------------------
% ------------------------------------------------------------------------------
% ------------------------------------------------------------------------------
% ------------------------------------------------------------------------------
% ------------------------------------------------------------------------------
% Pacotes ABNTEX

% citações padrão ABNT
\usepackage[alf]{abntex2cite}
% ------------------------------------------------------------------------------



% ------------------------------------------------------------------------------
% ------------------------------------------------------------------------------
% ------------------------------------------------------------------------------
% ------------------------------------------------------------------------------
% ------------------------------------------------------------------------------
% Macros do documento
\def\universidade{UNINASSAU}

\def\idcurso{99}
\def\curso{Nome do Curso ou Tema}

\def\iddisciplina{99}
\def\disciplina{Nome da Disciplina ou Descrição Detalhada do Tópico em Estudo}

\def\professor{João Ferreira da Silva Júnior}
\def\professorfoto{img/sec01/professor.png}
\def\professoremail{joaoferreirape@gmail.com}
\def\professorsite{https://joaoferreirape.wordpress.com}
\def\professorlattes{http://lattes.cnpq.br/8904695743376784}

\def\semana{}
\def\diadaaula{}
\def\mesdaaula{}
\def\anodaaula{}
\def\semestre{2023-1}
% \def\data{\diadaaula~de~\mesdaaula~de~\anodaaula}
\def\data{Semestre: \semestre}

%\title[\iddisciplina]{\textbf{\curso}}
%\title[\iddisciplina]{\textbf{\disciplina}}
\title[\curso]{\curso}

%\subtitle{\iddisciplina~---~\disciplina}
\subtitle[\disciplina]{\disciplina}

\subject{\semestre~---~\curso~---~\disciplina}

%\author[SILVA JUNIOR, J. F. $\therefore$]{\href{\professorlattes}{\textbf{\professor}}}
\author[SILVA JUNIOR, J. F. $\therefore$]{\href{\professorlattes}{\professor}}

\logo{\includegraphics[height=6ex]{img/logo/brasao.png}}

\institute[UNINASSAU]{\includegraphics[width=40ex]{img/logo/logo.png}}

\date{\data}
% ------------------------------------------------------------------------------



% ------------------------------------------------------------------------------
% ------------------------------------------------------------------------------
% ------------------------------------------------------------------------------
% ------------------------------------------------------------------------------
% ------------------------------------------------------------------------------
% configurações do pacote hyperref, precisa estar aqui devido as macros de dados do documento
\makeatletter
\hypersetup{% http://www.tug.org/applications/hyperref/manual.html#x1-120003.8
  anchorcolor={black},                      % set color of anchors
  % backref=false,                            % do bibliographical back references
  baseurl={\professorlattes},               % set base URL for document
  urlcolor={black},%{blue},                 % url color link
  % bookmarks={true},                         % make bookmarks
  bookmarksnumbered={false},                % put section numbers in bookmarks
  bookmarksopen={true},                     % open up bookmark tree
  bookmarksopenlevel={\maxdimen},           % level to which bookmarks are open
  % bookmarksdepth,                           % <no-value>, <number>, <part,chapter,section,...>
  bookmarkstype={toc},                      % to specify which ‘toc’ file to mimic
  % breaklinks={false},                       % allow links to break over lines
  % citebordercolor=0 1 0,                    % color of border around cites
  citecolor={black},%{green},               % color of citation links
  colorlinks={false},
  % false,                                    % color links
  % true,                                     % (tex4ht, dviwindo)
  % debug=false,                              % provide details of anchors defined; same as verbose
  % destlabel=false,                          % destinations are named by the first \label after the anchor creation
  draft={false},                            % do not do any hyperlinking
  % filebordercolor=0 .5 .5,                  % color of border around file links
  filecolor={black},%{cyan},                % color of file links
  final={true},                             % opposite of option draft
  % frenchlinks=false,                        % use small caps instead of color for links
  % hyperfigures=false,                       % make figures hyper links
  % hyperfootnotes=true,                      % set up hyperlinked footnotes
  % hyperindex=true,                          % set up hyperlinked indices
  % linkbordercolor=1 0 0,                    % color of border around links
  linkcolor={black},%{red},                 % color of links
  % linktocpage={false},                      % make page number, not text, be link on TOC, LOF and LOT
  % menubordercolor=1 0 0,                    % color of border around menu links
  menucolor={black},%{red},                 % color for menu links
  pageanchor={true},                        % put an anchor on every page
  % pagebackref={true},                       % backreference by page number
  pdfauthor={\professor~<\professoremail>}, % text for PDF Author field
  % pdfborder=0 0 1,                          % width of PDF link border
  % 0 0 0,                                    % (colorlinks)
  % pdfcenterwindow=true,                     % position the document window in the center of the screen
  pdfcreator={LaTeX with hyperref package, feet: \professor},
  % text for PDF Creator field
  pdfdisplaydoctitle=true,                  % display document title instead of file name in title bar
  pdfkeywords={{\universidade};{\curso};{\disciplina};{\professor}},
  % text for PDF Keywords field
  pdflang={pt-BR},                          % PDF language identifier (RFC 3066)
  pdfproducer={\professor~<\professoremail>},   % text for PDF Producer field
  % pdfstartpage=1,                           % page at which PDF document opens
  pdfsubject={\semestre~---~\curso~---~\disciplina},  % text for PDF Subject field
  pdftitle={\curso~---~\disciplina},    % text for PDF Title field
}
\makeatother
% ------------------------------------------------------------------------------



% ------------------------------------------------------------------------------
% ------------------------------------------------------------------------------
% ------------------------------------------------------------------------------
% ------------------------------------------------------------------------------
% ------------------------------------------------------------------------------
\begin{document}



% ------------------------------------------------------------------------------
% ------------------------------------------------------------------------------
% ------------------------------------------------------------------------------
% ------------------------------------------------------------------------------
% ------------------------------------------------------------------------------
% ------------------------------------------------------------------------------
% ------------------------------------------------------------------------------
% ------------------------------------------------------------------------------
% ------------------------------------------------------------------------------
% ------------------------------------------------------------------------------
\section*{Capa Externa}\label{sec:capa-externa}



{\usebackgroundtemplate{\includegraphics[width=\paperwidth,height=\paperheight,keepaspectratio]{img/logo/background-amarelo.png}}
  \begin{frame}[plain]

    ~

  \end{frame}
} % usebackgroundtemplate



% ------------------------------------------------------------------------------
% ------------------------------------------------------------------------------
% ------------------------------------------------------------------------------
% ------------------------------------------------------------------------------
% ------------------------------------------------------------------------------
% ------------------------------------------------------------------------------
% ------------------------------------------------------------------------------
% ------------------------------------------------------------------------------
% ------------------------------------------------------------------------------
% ------------------------------------------------------------------------------
\section*{Capa}\label{sec:capa}



{\usebackgroundtemplate{{\transparent{0.1}\includegraphics[width=\paperwidth,height=\paperheight,keepaspectratio]{img/logo/background-amarelo.png}}}
  \begin{frame}[plain,t]

    \titlepage

  \end{frame}
} % usebackgroundtemplate



% ------------------------------------------------------------------------------
% ------------------------------------------------------------------------------
% ------------------------------------------------------------------------------
% ------------------------------------------------------------------------------
% ------------------------------------------------------------------------------
% ------------------------------------------------------------------------------
% ------------------------------------------------------------------------------
% ------------------------------------------------------------------------------
% ------------------------------------------------------------------------------
% ------------------------------------------------------------------------------
\section*{Agenda}\label{sec:agenda}



\begin{frame}[allowframebreaks,t]\frametitle{Agenda ---}
%\begin{frame}[allowframebreaks,t]
%\begin{frame}[t]\frametitle{Agenda}

  \tableofcontents{}
  % \tableofcontents[sectionstyle=show/hide/shaded,subsectionstyle=show/hide/shaded]
  % \tableofcontents[sectionstyle=show,subsectionstyle=show]
  % \tableofcontents[hideallsubsections]
  % \tableofcontents[currentsection,currentsubsection]

\end{frame}



% ------------------------------------------------------------------------------
% ------------------------------------------------------------------------------
% ------------------------------------------------------------------------------
% ------------------------------------------------------------------------------
% ------------------------------------------------------------------------------
% ------------------------------------------------------------------------------
% ------------------------------------------------------------------------------
% ------------------------------------------------------------------------------
% ------------------------------------------------------------------------------
% ------------------------------------------------------------------------------
% !TeX encoding = UTF-8
% !TeX spellcheck = pt_BR
% !TeX root = 20232-CODDISC-JOAO-FERREIRA.tex
% ------------------------------------------------------------------------------

% ------------------------------------------------------------------------------
% Template para aulas da UNINASSAU
% Autor: João Ferreira da Silva Júnior
% Versão: 0.4.4 - 2023-07-19
% ------------------------------------------------------------------------------



% ------------------------------------------------------------------------------
% ------------------------------------------------------------------------------
% ------------------------------------------------------------------------------
% ------------------------------------------------------------------------------
% ------------------------------------------------------------------------------
% ------------------------------------------------------------------------------
% ------------------------------------------------------------------------------
% ------------------------------------------------------------------------------
% ------------------------------------------------------------------------------
% ------------------------------------------------------------------------------
\section*{Professor}\label{sec:professor}



% \begin{frame}[allowframebreaks,t]\frametitle{Professor ---}
% \begin{frame}[t]\frametitle{Professor}
\begin{frame}[t]

    \begin{columns}[onlytextwidth,T]

        \begin{column}[T]{0.49\linewidth}
            \begin{figure}[htb]
              \centering{}
              \includegraphics[width=0.5\linewidth,height=0.3\textheight,keepaspectratio]{\professorfoto}

              \textbf{\href{\professorlattes}{\professor}}
            \end{figure}

            \begin{itemize}
              \justifying{}
              \setlength\itemsep{1em}
              \item \href{\professorsite}{\professorsite}
              \item \href{\professorlattes}{\professorlattes}
              \item Formado em Análise e Desenvolvimento de Sistemas em 2013 (UNOPAR).
            \end{itemize}
        \end{column}

        \begin{column}{0.5\linewidth}
            \small
            \begin{itemize}\small
                \justifying{}
                % \setlength\itemsep{1em}
                \item Pós-graduado em nível de especialista em Cloud Computing, 2021 (UPE).
                \item Pós-graduado em nível de especialista em I.A., 2022 (UPE).
                \item Mestre em Engenharia Elétrica, título obtido em 2017 (UFPE) com o tema Detecção de Perdas em Sistemas de Distribuição de Água.
                \item Doutorando em Engenharia Elétrica (UFPE).
                \item Professor da UNINASSAU nos cursos de Análise e Desenvolvimento de Sistemas, Ciência da Computação, Engenharia da Computação, Redes de Computadores, e Sistemas de Informação.
            \end{itemize}
        \end{column}

    \end{columns}

\end{frame}


% !TeX encoding = UTF-8
% !TeX spellcheck = pt_BR
% !TeX root = 20231-CODDISC-JOAO-FERREIRA.tex
% ------------------------------------------------------------------------------

% ------------------------------------------------------------------------------
% Template para aulas da UNINASSAU
% Autor: João Ferreira da Silva Júnior
% Versão: 0.4.3 - 2023-02-04
% ------------------------------------------------------------------------------



% ------------------------------------------------------------------------------
% ------------------------------------------------------------------------------
% ------------------------------------------------------------------------------
% ------------------------------------------------------------------------------
% ------------------------------------------------------------------------------
% ------------------------------------------------------------------------------
% ------------------------------------------------------------------------------
% ------------------------------------------------------------------------------
% ------------------------------------------------------------------------------
% ------------------------------------------------------------------------------
\section[Plano de Ensino]{Plano de Ensino}\label{sec:plano-ensino}



{\usebackgroundtemplate{{\transparent{0.1}\includegraphics[width=\paperwidth,height=\paperheight,keepaspectratio]{img/logo/background-branco.png}}}
% \begin{frame}[allowframebreaks,t]%\frametitle{ ---}
  \begin{frame}[plain]%\frametitle{}

    \vfill

    \centering{}
    \Huge{\textbf{Plano de Ensino}}

    \vfill

\end{frame}
} % usebackgroundtemplate



% ------------------------------------------------------------------------------
% ------------------------------------------------------------------------------
% ------------------------------------------------------------------------------
% ------------------------------------------------------------------------------
% ------------------------------------------------------------------------------
\subsection[Ementa]{Ementa}\label{subsec:plano-ensino-ementa}



% \begin{frame}[allowframebreaks,t]\frametitle{Ementa ---}
\begin{frame}[t]\frametitle{Ementa}

    \begin{block}{Objetivo}
      \justifying{}
      Analisar, avaliar, compreender, conceituar, conhecer, descrever, desenvolver, documentar, empregar, identificar, sintetizar, $\ldots$
      assuntos importantes acerca do tema ou tópico, através do estudo teórico e desenvolvimento prático.
    \end{block}

    \begin{block}{Competências}
      \justifying{}
      \begin{itemize}
        \justifying{}
%        \setlength\itemsep{1em}
        \item Competência 1.
        \item Competência 2.
        \item Competência 3.
        \item Competência 4.
        \item Competência 5.
      \end{itemize}
    \end{block}

\end{frame}



% ------------------------------------------------------------------------------
% ------------------------------------------------------------------------------
% ------------------------------------------------------------------------------
% ------------------------------------------------------------------------------
% ------------------------------------------------------------------------------
\subsection[Conteúdo Programático]{Conteúdo Programático}\label{subsec:plano-ensino-conteudo}



% \begin{frame}[allowframebreaks,t]\frametitle{Conteúdo Programático ---}
\begin{frame}[t]\frametitle{Conteúdo Programático}

  \begin{columns}[onlytextwidth,T]

    \begin{column}{0.5\linewidth}
      \begin{block}{Unidade I:}
        \begin{itemize}\small
          \justifying{}
          % \setlength\itemsep{1em}
          \item Item 1.
          \begin{itemize}\scriptsize
            \justifying{}
            \item Subitem 1.1.
            \item Subitem 1.2.
          \end{itemize}
          \item Item 2.
          \item Item 3.
          \item Item 4.
        \end{itemize}
      \end{block}
    \end{column}

    \begin{column}{0.5\linewidth}
      \begin{block}{Unidade II:}
        \begin{itemize}\small
          \justifying{}
          % \setlength\itemsep{1em}
          \item Item 5.
          \item Item 6.
          \item Item 7.
          \item Item 8.
        \end{itemize}
      \end{block}
    \end{column}

  \end{columns}

\end{frame}



% ------------------------------------------------------------------------------
% ------------------------------------------------------------------------------
% ------------------------------------------------------------------------------
% ------------------------------------------------------------------------------
% ------------------------------------------------------------------------------
\subsection[Metodologia de Ensino]{Metodologia de Ensino}\label{subsec:plano-ensino-metodologia}



% \begin{frame}[allowframebreaks,t]\frametitle{Metodologia de Ensino ---}
\begin{frame}[t]\frametitle{Metodologia de Ensino}

  \begin{block}{}
    \justifying{}
    \Large
    A disciplina, dependendo de sua natureza, pode ser ministrada através de conteúdos teóricos, conteúdos práticos, aulas de campo em instituições específicas e ainda pode utilizar recursos de exposições dialogadas, grupos de discussão, seminários, debates competitivos, apresentação e discussão de filmes e casos práticos, onde os conteúdos podem ser trabalhados mais dinamicamente, estimulando o senso crítico e científico dos alunos.
  \end{block}

\end{frame}



% ------------------------------------------------------------------------------
% ------------------------------------------------------------------------------
% ------------------------------------------------------------------------------
% ------------------------------------------------------------------------------
% ------------------------------------------------------------------------------
\subsection[Metodologia de Avaliação]{Metodologia de Avaliação}\label{subsec:plano-ensino-avaliacao}



% \begin{frame}[allowframebreaks,t]\frametitle{Metodologia de Avaliação ---}
\begin{frame}[t]\frametitle{Metodologia de Avaliação}

  \begin{block}{}
    \justifying{}
    \Large
    No decorrer de cada período letivo são desenvolvidas 02 (duas) avaliações por disciplina, para efeito do cálculo da média parcial. A média parcial é calculada pela média aritmética das duas avaliações efetuadas. O aluno que alcançar a média parcial maior ou igual a 7,0 (sete) é considerado aprovado. O aluno que não alcançar a média parcial faz em exame final onde precisa alcançar média final maior ou igual a 5,0. São aplicadas avaliações dos tipos: provas teóricas, provas práticas, seminários, trabalhos individuais ou em grupo e outras atividades em classe e extraclasse. O exame final é, obrigatoriamente, prova escrita.
  \end{block}

\end{frame}



% ------------------------------------------------------------------------------
% ------------------------------------------------------------------------------
% ------------------------------------------------------------------------------
% ------------------------------------------------------------------------------
% ------------------------------------------------------------------------------
\subsection[Bibliografia Sugerida]{Bibliografia Sugerida}\label{subsec:plano-ensino-bibliografia}



% \begin{frame}[allowframebreaks,t]\frametitle{Bibliografia Básica ---}
\begin{frame}[t]\frametitle{Bibliografia Básica}

  \begin{itemize}
    \justifying{}
    \setlength\itemsep{1em}
    \item Exemplos de citações.
    \item Redes de Computadores. 5. ed. [S.l.]: Pearson Prentice Hall. \cite{Tanenbaum2011}
    \item Sistemas Operacionais Modernos. 4. ed. [S.l.]: Pearson Prentice Hall. \cite{Tanenbaum2016}
  \end{itemize}

\end{frame}



% \begin{frame}[allowframebreaks,t]\frametitle{Bibliografia Complementar ---}
\begin{frame}[t]\frametitle{Bibliografia Complementar}

  \begin{itemize}
    \justifying{}
    \setlength\itemsep{1em}
    \item Exemplos de citações.
    \item Redes de Computadores. 5. ed. [S.l.]: Pearson Prentice Hall. \cite{Tanenbaum2011}
    \item Sistemas Operacionais Modernos. 4. ed. [S.l.]: Pearson Prentice Hall. \cite{Tanenbaum2016}
  \end{itemize}

\end{frame}


% !TeX encoding = UTF-8
% !TeX spellcheck = pt_BR
% !TeX root = 20231-CODDISC-JOAO-FERREIRA.tex
% ------------------------------------------------------------------------------

% ------------------------------------------------------------------------------
% Template para aulas da UNINASSAU
% Autor: João Ferreira da Silva Júnior
% Versão: 0.4.4 - 2023-07-19
% ------------------------------------------------------------------------------



% ------------------------------------------------------------------------------
% ------------------------------------------------------------------------------
% ------------------------------------------------------------------------------
% ------------------------------------------------------------------------------
% ------------------------------------------------------------------------------
% ------------------------------------------------------------------------------
% ------------------------------------------------------------------------------
% ------------------------------------------------------------------------------
% ------------------------------------------------------------------------------
% ------------------------------------------------------------------------------
\section[Planejamento da Disciplina]{Planejamento da Disciplina}\label{sec:planejamento}



{\usebackgroundtemplate{{\transparent{0.1}\includegraphics[width=\paperwidth,height=\paperheight,keepaspectratio]{img/logo/background-branco.png}}}
  % \begin{frame}[allowframebreaks,t]%\frametitle{ ---}
  \begin{frame}[plain]%\frametitle{}

    \vfill

    \centering{}
    \Huge{\textbf{Planejamento da Disciplina}}

    \vfill

  \end{frame}
} % usebackgroundtemplate



% ------------------------------------------------------------------------------
% ------------------------------------------------------------------------------
% ------------------------------------------------------------------------------
% ------------------------------------------------------------------------------
% ------------------------------------------------------------------------------
\subsection[Calendário Acadêmico]{Calendário Acadêmico}\label{subsec:planejamento-calendario}



% \begin{frame}[allowframebreaks,t]\frametitle{Calendário Acadêmico ---}
\begin{frame}[t]\frametitle{Calendário Acadêmico}

  \begin{itemize}
    \justifying{}
    \setlength\itemsep{1em}
    \item 01/01/2001 a 01/01/2001 --- início das aulas.
    \item 02/02/2002 a 02/02/2002 --- lista de exercícios, seminários e projetos.
    \item 02/02/2002 a 02/02/2002 --- 1ª avaliação.
    \item 03/03/2003 a 03/03/2003 --- lista de exercícios, seminários e projetos.
    \item 03/03/2003 a 03/03/2003 --- 2ª avaliação.
    \item 04/04/2004 a 04/04/2004 --- 2ª chamada.
    \item 05/05/2005 a 05/05/2005 --- avaliação final.
    \item O detalhamento e datas específicas constam no portal do aluno.
  \end{itemize}

\end{frame}



% ------------------------------------------------------------------------------
% ------------------------------------------------------------------------------
% ------------------------------------------------------------------------------
% ------------------------------------------------------------------------------
% ------------------------------------------------------------------------------
\subsection[Participação]{Participação}\label{subsec:planejamento-participacao}



% \begin{frame}[allowframebreaks,t]\frametitle{Participação ---}
\begin{frame}[t]\frametitle{Participação}

  \begin{itemize}
    \justifying{}
    \setlength\itemsep{1em}
    \item Frequência mínima de 75\% das aulas, determinação institucional.
    \item Participação em aula e interação útil entre os alunos.
    \item Não às conversas que causam dispersão.
    \item Questionamentos pertinentes e ajuda aos demais colegas de turma.
    \item Adequação ao conteúdo ministrado e contexto da disciplina.
  \end{itemize}

\end{frame}



% ------------------------------------------------------------------------------
% ------------------------------------------------------------------------------
% ------------------------------------------------------------------------------
% ------------------------------------------------------------------------------
% ------------------------------------------------------------------------------
\subsection[Lista de Exercícioss]{Lista de Exercícioss}\label{subsec:planejamento-lista}



% \begin{frame}[allowframebreaks,t]\frametitle{Lista de Exercícioss ---}
\begin{frame}[t]\frametitle{Lista de Exercícioss}

  \begin{itemize}
    \justifying{}
    \setlength\itemsep{1em}
    \item Poderá, ou não, haver listas de exercícios, que irão valer ponto e servirão como base de estudo, mas não de todo o conteúdo ministrado, para as avaliações.
    \item A eventual pontuação por lista de exercício será avaliada e computada individualmente.
    \item A(s) lista(s) de exercícios aplicadas até a primeira avaliação, terá(ão) conteúdo relacionado e contará eventual pontuação para esta.
    \item A(s) lista(s) de exercícios apicadas após a primeira avaliação e até a segunda, terá(ão) conteúdo relacionado e contará eventual pontuação para esta.
    \item Nenhuma eventual pontuação valerá para a segunda chamada ou avaliação final.
    \item A(s) lista(s) de exercício valerá(ão) até 10\% (dez por cento) da respectiva avaliação.
  \end{itemize}

\end{frame}



% ------------------------------------------------------------------------------
% ------------------------------------------------------------------------------
% ------------------------------------------------------------------------------
% ------------------------------------------------------------------------------
% ------------------------------------------------------------------------------
\subsection[Seminários]{Seminários}\label{subsec:planejamento-seminarios}



% \begin{frame}[allowframebreaks,t]\frametitle{Seminários ---}
\begin{frame}[t]\frametitle{Seminários}

  \begin{itemize}
    \justifying{}
    \setlength\itemsep{1em}
    \item Poderá, ou não, haver apresentação de seminários, que irão valer ponto e servirão como base de estudo, mas não de todo conteúdo ministrado, para as avaliações.
    \item A eventual pontuação por seminário será avaliada e computada individualmente.
    \item Poderá ser feito individualmente ou em grupo de até 5 (cinco) membros, mas atentem de que aleatoriamente serão feitos questionamentos a qualquer um dos membros, e, na incapacidade de resposta, a falha será atribuída a todo o grupo.
    \item O seminário deverá ser apresentado para toda a turma quando serão observados e avaliados os aspectos do tempo de apresentação; conteúdo; clareza; domínio; preparo geral e recursos utilizados.
    \item A nota atribuída valerá até 10\% (dez por cento) da respectiva avaliação.
  \end{itemize}

\end{frame}



% ------------------------------------------------------------------------------
% ------------------------------------------------------------------------------
% ------------------------------------------------------------------------------
% ------------------------------------------------------------------------------
% ------------------------------------------------------------------------------
\subsection[Projeto]{Projeto}\label{subsec:planejamento-projeto}



% \begin{frame}[allowframebreaks,t]\frametitle{Projeto ---}
\begin{frame}[t]\frametitle{Projeto}

  \begin{itemize}
    \justifying{}
    \setlength\itemsep{1em}
    \item Poderá, ou não, haver projeto, voltado à aplicação dos conceitos aprendidos em sala de aula, nas listas de exercícios, nos seminários e na bibliografia que consta no Plano de Ensino.
    \item A eventual pontuação por projeto será avaliada e computada individualmente.
    \item Poderá ser feito individualmente ou em grupo de até 5 (cinco) membros, mas atentem de que aleatoriamente serão feitos questionamentos a qualquer um dos membros, e, na incapacidade de resposta, a falha será atribuída a todo o grupo.
    \item O projeto será avaliado em apresentação para toda a turma, quando deverão ser demonstrados todos os processos desenvolvidos na pesquisa e elaboração, bem como o eventual funcionamento do projeto.
    \item A nota atribuída valerá até 20\% (vinte por cento) da respectiva avaliação.
  \end{itemize}

\end{frame}



% ------------------------------------------------------------------------------
% ------------------------------------------------------------------------------
% ------------------------------------------------------------------------------
% ------------------------------------------------------------------------------
% ------------------------------------------------------------------------------
\subsection[Avaliações]{Avaliações}\label{subsec:planejamento-avaliacoes}



% \begin{frame}[allowframebreaks,t]\frametitle{Avaliações ---}
\begin{frame}[t]\frametitle{Avaliações}

  \begin{columns}[onlytextwidth,T]

    \begin{column}{0.48\linewidth}
      \begin{block}{1ª e 2ª avaliações:}
        \begin{itemize}
            \justifying{}
            %\setlength\itemsep{1em}
            \item Conteúdo estudado em cada Unidade:
            \begin{itemize}
              \justifying{}
              \item Bibliografia.
              \item Listas de exercícios.
              \item Seminários que foram apresentados.
              \item Projetos que foram propostos.
            \end{itemize}
            \item Cálculo da nota:
            \begin{itemize}
              \justifying{}
              \item Listas de exercícios até +10\%, se houver.
              \item Seminários até +10\%, se houver.
              \item Projetos até +20\%, se houver.
              \item Avaliação até 60\%, com todos os extras.
              \item Avaliação até 100\%, sem os extras.
            \end{itemize}
        \end{itemize}
      \end{block}
    \end{column}

    \begin{column}{0.48\linewidth}
        \begin{block}{2ª chamada e avaliação final:}
          \begin{itemize}
            \justifying{}
            %\setlength\itemsep{1em}
            \item Conteúdo estudado em ambas Unidades:
            \begin{itemize}
              \justifying{}
              \item Bibliografia.
              \item Listas de exercícios.
              \item Seminários que foram apresentados.
              \item Projetos que foram propostos.
            \end{itemize}
            \item Cálculo da nota:
            \begin{itemize}
              \justifying{}
              \item Não haverá atividade extra.
              \item Qualquer ponto de atividade extra, feita anteriormente, não será computado para as notas da 2ª chamada ou avaliação final.
              \item A nota de cada avaliação valerá 100\%.
            \end{itemize}
          \end{itemize}
        \end{block}
    \end{column}

  \end{columns}

\end{frame}



% ------------------------------------------------------------------------------
% ------------------------------------------------------------------------------
% ------------------------------------------------------------------------------
% ------------------------------------------------------------------------------
% ------------------------------------------------------------------------------
\subsection[Comunicação da Turma]{Comunicação da Turma}\label{subsec:planejamento-comunicacao-turma}



% \begin{frame}[allowframebreaks,t]\frametitle{Comunicação da Turma ---}
\begin{frame}[t]\frametitle{Comunicação da Turma}

  \begin{figure}[htb]
    \centering{}
    \includegraphics[width=0.07\linewidth]{img/sec03/telegram.png}

    \textbf{\href{https://telegram.org}{Telegram}}
  \end{figure}

  \begin{columns}[onlytextwidth,T]

    \begin{column}{0.48\linewidth}
      \begin{block}{Vantagens:}
        \begin{itemize}
          \justifying{}
          % \setlength\itemsep{1em}
          \item Grupo pode ter muitos membros e a troca de arquivos é bem otimizada.
          \item Histórico permanente, quem entra depois poderá ver todas as mensagens já trocadas no grupo.
          \item Ótimo aplicativo para dispositivos móveis, desktop e web.
        \end{itemize}
      \end{block}
    \end{column}

    \begin{column}{0.48\linewidth}
      \begin{block}{Disponível para:}
        \begin{itemize}
          \justifying{}
          % \setlength\itemsep{1em}
          \item \href{https://telegram.org/dl/android}{Android}
          \item \href{https://telegram.org/dl/ios}{iPhone/iPad}
          \item \href{https://desktop.telegram.org}{PC/Mac/Linux}
          \item \href{https://web.telegram.org}{Web}
          \item \href{https://macos.telegram.org}{macOS}
          \item \href{https://web.telegram.org}{https://web.telegram.org}
        \end{itemize}
      \end{block}
    \end{column}

  \end{columns}

\end{frame}


% !TeX encoding = UTF-8
% !TeX spellcheck = pt_BR
% !TeX root = 20231-CODDISC-JOAO-FERREIRA.tex
% ------------------------------------------------------------------------------

% ------------------------------------------------------------------------------
% Template para aulas da UNINASSAU
% Autor: João Ferreira da Silva Júnior
% Versão: 0.4.4 - 2023-07-19
% ------------------------------------------------------------------------------



% ------------------------------------------------------------------------------
% ------------------------------------------------------------------------------
% ------------------------------------------------------------------------------
% ------------------------------------------------------------------------------
% ------------------------------------------------------------------------------
% ------------------------------------------------------------------------------
% ------------------------------------------------------------------------------
% ------------------------------------------------------------------------------
% ------------------------------------------------------------------------------
% ------------------------------------------------------------------------------
\section[Introdução]{Introdução}\label{sec:introducao}



{\usebackgroundtemplate{{\transparent{0.1}\includegraphics[width=\paperwidth,height=\paperheight,keepaspectratio]{img/logo/background-branco.png}}}
  % \begin{frame}[allowframebreaks,t]%\frametitle{ ---}
  \begin{frame}[plain]%\frametitle{}

    \vfill

    \centering{}
    \Huge{\textbf{Introdução}}

    \vfill

  \end{frame}
} % usebackgroundtemplate



% ------------------------------------------------------------------------------
% ------------------------------------------------------------------------------
% ------------------------------------------------------------------------------
% ------------------------------------------------------------------------------
% ------------------------------------------------------------------------------
\subsection[Contextualização]{Contextualização}\label{subsec:introducao-contextualizacao}



% \begin{frame}[allowframebreaks,t]\frametitle{Contextualização ---}
\begin{frame}[t]\frametitle{Contextualização}

  \begin{itemize}
    \justifying{}
    \setlength\itemsep{1em}
    \item Item 1.
    \item Item 2.
    \item Item 3.
  \end{itemize}

\end{frame}



% ------------------------------------------------------------------------------
% ------------------------------------------------------------------------------
% ------------------------------------------------------------------------------
% ------------------------------------------------------------------------------
% ------------------------------------------------------------------------------
\subsection[Exercícios]{Exercícios}\label{subsec:introducao-exercicios}



% \begin{frame}[allowframebreaks,t]\frametitle{Exercícios ---}
\begin{frame}[t]\frametitle{Exercícios}

  \begin{itemize}
    \justifying{}
    \setlength\itemsep{1em}
    \item Questão 1.
    \item Questão 2.
    \item Questão 3.
  \end{itemize}

\end{frame}


% !TeX encoding = UTF-8
% !TeX spellcheck = pt_BR
% !TeX root = 20231-CODDISC-JOAO-FERREIRA.tex
% ------------------------------------------------------------------------------

% ------------------------------------------------------------------------------
% Template para aulas da UNINASSAU
% Autor: João Ferreira da Silva Júnior
% Versão: 0.4.3 - 2023-02-04
% ------------------------------------------------------------------------------



% ------------------------------------------------------------------------------
% ------------------------------------------------------------------------------
% ------------------------------------------------------------------------------
% ------------------------------------------------------------------------------
% ------------------------------------------------------------------------------
% ------------------------------------------------------------------------------
% ------------------------------------------------------------------------------
% ------------------------------------------------------------------------------
% ------------------------------------------------------------------------------
% ------------------------------------------------------------------------------
\section[Tema 01]{Tema 01}\label{sec:tema01}



{\usebackgroundtemplate{{\transparent{0.1}\includegraphics[width=\paperwidth,height=\paperheight,keepaspectratio]{img/logo/background-branco.png}}}
  % \begin{frame}[allowframebreaks,t]%\frametitle{ ---}
  \begin{frame}[plain]%\frametitle{}

    \vfill

    \centering{}
    \Huge{\textbf{Tema 01}}

    \vfill

  \end{frame}
} % usebackgroundtemplate



% ------------------------------------------------------------------------------
% ------------------------------------------------------------------------------
% ------------------------------------------------------------------------------
% ------------------------------------------------------------------------------
% ------------------------------------------------------------------------------
\subsection[Tópico 01]{Tópico 01}\label{subsec:tema01-topico01}



% \begin{frame}[allowframebreaks,t]\frametitle{Tópico 01 ---}
\begin{frame}[t]\frametitle{Tópico 01}

  \begin{itemize}
    \justifying{}
    \setlength\itemsep{1em}
    \item Item 1.
    \item Item 2.
    \item Item 3.
  \end{itemize}

\end{frame}



% ------------------------------------------------------------------------------
% ------------------------------------------------------------------------------
% ------------------------------------------------------------------------------
% ------------------------------------------------------------------------------
% ------------------------------------------------------------------------------
\subsection[Tópico 02]{Tópico 02}\label{subsec:tema01-topico02}



% \begin{frame}[allowframebreaks,t]\frametitle{Tópico 02 ---}
\begin{frame}[t]\frametitle{Tópico 02}
  
  \begin{itemize}
    \justifying{}
    \setlength\itemsep{1em}
    \item Item 1.
    \item Item 2.
    \item Item 3.
  \end{itemize}
  
\end{frame}



% ------------------------------------------------------------------------------
% ------------------------------------------------------------------------------
% ------------------------------------------------------------------------------
% ------------------------------------------------------------------------------
% ------------------------------------------------------------------------------
\subsection[Exercícios]{Exercícios}\label{subsec:tema01-exercicios}



% \begin{frame}[allowframebreaks,t]\frametitle{Exercícios ---}
\begin{frame}[t]\frametitle{Exercícios}

\begin{itemize}
  \justifying{}
  \setlength\itemsep{1em}
  \item Questão 1.
  \item Questão 2.
  \item Questão 3.
\end{itemize}

\end{frame}


% !TeX encoding = UTF-8
% !TeX spellcheck = pt_BR
% !TeX root = 20231-CODDISC-JOAO-FERREIRA.tex
% ------------------------------------------------------------------------------

% ------------------------------------------------------------------------------
% Template para aulas da UNINASSAU
% Autor: João Ferreira da Silva Júnior
% Versão: 0.4.4 - 2023-07-19
% ------------------------------------------------------------------------------



% ------------------------------------------------------------------------------
% ------------------------------------------------------------------------------
% ------------------------------------------------------------------------------
% ------------------------------------------------------------------------------
% ------------------------------------------------------------------------------
% ------------------------------------------------------------------------------
% ------------------------------------------------------------------------------
% ------------------------------------------------------------------------------
% ------------------------------------------------------------------------------
% ------------------------------------------------------------------------------
\section[Tema 02]{Tema 02}\label{sec:tema02}



{\usebackgroundtemplate{{\transparent{0.1}\includegraphics[width=\paperwidth,height=\paperheight,keepaspectratio]{img/logo/background-branco.png}}}
  % \begin{frame}[allowframebreaks,t]%\frametitle{ ---}
  \begin{frame}[plain]%\frametitle{}

    \vfill

    \centering{}
    \Huge{\textbf{Tema 02}}

    \vfill

  \end{frame}
} % usebackgroundtemplate



% ------------------------------------------------------------------------------
% ------------------------------------------------------------------------------
% ------------------------------------------------------------------------------
% ------------------------------------------------------------------------------
% ------------------------------------------------------------------------------
\subsection[Tópico 01]{Tópico 01}\label{subsec:tema02-topico01}



% \begin{frame}[allowframebreaks,t]\frametitle{Tópico 01 ---}
\begin{frame}[t]\frametitle{Tópico 01}

  \begin{itemize}
    \justifying{}
    \setlength\itemsep{1em}
    \item Item 1.
    \item Item 2.
    \item Item 3.
  \end{itemize}

\end{frame}



% ------------------------------------------------------------------------------
% ------------------------------------------------------------------------------
% ------------------------------------------------------------------------------
% ------------------------------------------------------------------------------
% ------------------------------------------------------------------------------
\subsection[Tópico 02]{Tópico 02}\label{subsec:tema02-topico02}



% \begin{frame}[allowframebreaks,t]\frametitle{Tópico 02 ---}
\begin{frame}[t]\frametitle{Tópico 02}

  \begin{itemize}
    \justifying{}
    \setlength\itemsep{1em}
    \item Item 1.
    \item Item 2.
    \item Item 3.
  \end{itemize}

\end{frame}



% ------------------------------------------------------------------------------
% ------------------------------------------------------------------------------
% ------------------------------------------------------------------------------
% ------------------------------------------------------------------------------
% ------------------------------------------------------------------------------
\subsection[Exercícios]{Exercícios}\label{subsec:tema02-exercicios}



% \begin{frame}[allowframebreaks,t]\frametitle{Exercícios ---}
\begin{frame}[t]\frametitle{Exercícios}

\begin{itemize}
  \justifying{}
  \setlength\itemsep{1em}
  \item Questão 1.
  \item Questão 2.
  \item Questão 3.
\end{itemize}

\end{frame}


% !TeX encoding = UTF-8
% !TeX spellcheck = pt_BR
% !TeX root = 20232-CODDISC-JOAO-FERREIRA.tex
% ------------------------------------------------------------------------------

% ------------------------------------------------------------------------------
% Template para aulas da UNINASSAU
% Autor: João Ferreira da Silva Júnior
% Versão: 0.4.4 - 2023-07-19
% ------------------------------------------------------------------------------



% ------------------------------------------------------------------------------
% ------------------------------------------------------------------------------
% ------------------------------------------------------------------------------
% ------------------------------------------------------------------------------
% ------------------------------------------------------------------------------
% ------------------------------------------------------------------------------
% ------------------------------------------------------------------------------
% ------------------------------------------------------------------------------
% ------------------------------------------------------------------------------
% ------------------------------------------------------------------------------
\section[Tema 03]{Tema 03}\label{sec:tema03}



{\usebackgroundtemplate{{\transparent{0.1}\includegraphics[width=\paperwidth,height=\paperheight,keepaspectratio]{img/logo/background-branco.png}}}
  % \begin{frame}[allowframebreaks,t]%\frametitle{ ---}
  \begin{frame}[plain]%\frametitle{}

    \vfill

    \centering{}
    \Huge{\textbf{Tema 03}}

    \vfill

  \end{frame}
} % usebackgroundtemplate



% ------------------------------------------------------------------------------
% ------------------------------------------------------------------------------
% ------------------------------------------------------------------------------
% ------------------------------------------------------------------------------
% ------------------------------------------------------------------------------
\subsection[Tópico 01]{Tópico 01}\label{subsec:tema03-topico01}



% \begin{frame}[allowframebreaks,t]\frametitle{Tópico 01 ---}
\begin{frame}[t]\frametitle{Tópico 01}

  \begin{itemize}
    \justifying{}
    \setlength\itemsep{1em}
    \item Item 1.
    \item Item 2.
    \item Item 3.
  \end{itemize}

\end{frame}



% ------------------------------------------------------------------------------
% ------------------------------------------------------------------------------
% ------------------------------------------------------------------------------
% ------------------------------------------------------------------------------
% ------------------------------------------------------------------------------
\subsection[Tópico 02]{Tópico 02}\label{subsec:tema03-topico02}



% \begin{frame}[allowframebreaks,t]\frametitle{Tópico 02 ---}
\begin{frame}[t]\frametitle{Tópico 02}

  \begin{itemize}
    \justifying{}
    \setlength\itemsep{1em}
    \item Item 1.
    \item Item 2.
    \item Item 3.
  \end{itemize}

\end{frame}



% ------------------------------------------------------------------------------
% ------------------------------------------------------------------------------
% ------------------------------------------------------------------------------
% ------------------------------------------------------------------------------
% ------------------------------------------------------------------------------
\subsection[Exercícios]{Exercícios}\label{subsec:tema03-exercicios}



% \begin{frame}[allowframebreaks,t]\frametitle{Exercícios ---}
\begin{frame}[t]\frametitle{Exercícios}

\begin{itemize}
  \justifying{}
  \setlength\itemsep{1em}
  \item Questão 1.
  \item Questão 2.
  \item Questão 3.
\end{itemize}

\end{frame}


% !TeX encoding = UTF-8
% !TeX spellcheck = pt_BR
% !TeX root = 20232-CODDISC-JOAO-FERREIRA.tex
% ------------------------------------------------------------------------------

% ------------------------------------------------------------------------------
% Template para aulas da UNINASSAU
% Autor: João Ferreira da Silva Júnior
% Versão: 0.4.4 - 2023-07-19
% ------------------------------------------------------------------------------



% ------------------------------------------------------------------------------
% ------------------------------------------------------------------------------
% ------------------------------------------------------------------------------
% ------------------------------------------------------------------------------
% ------------------------------------------------------------------------------
% ------------------------------------------------------------------------------
% ------------------------------------------------------------------------------
% ------------------------------------------------------------------------------
% ------------------------------------------------------------------------------
% ------------------------------------------------------------------------------
\section[Tema 04]{Tema 04}\label{sec:tema04}



{\usebackgroundtemplate{{\transparent{0.1}\includegraphics[width=\paperwidth,height=\paperheight,keepaspectratio]{img/logo/background-branco.png}}}
  % \begin{frame}[allowframebreaks,t]%\frametitle{ ---}
  \begin{frame}[plain]%\frametitle{}

    \vfill

    \centering{}
    \Huge{\textbf{Tema 04}}

    \vfill

  \end{frame}
} % usebackgroundtemplate



% ------------------------------------------------------------------------------
% ------------------------------------------------------------------------------
% ------------------------------------------------------------------------------
% ------------------------------------------------------------------------------
% ------------------------------------------------------------------------------
\subsection[Tópico 01]{Tópico 01}\label{subsec:tema04-topico01}



% \begin{frame}[allowframebreaks,t]\frametitle{Tópico 01 ---}
\begin{frame}[t]\frametitle{Tópico 01}

  \begin{itemize}
    \justifying{}
    \setlength\itemsep{1em}
    \item Item 1.
    \item Item 2.
    \item Item 3.
  \end{itemize}

\end{frame}



% ------------------------------------------------------------------------------
% ------------------------------------------------------------------------------
% ------------------------------------------------------------------------------
% ------------------------------------------------------------------------------
% ------------------------------------------------------------------------------
\subsection[Tópico 02]{Tópico 02}\label{subsec:tema04-topico02}



% \begin{frame}[allowframebreaks,t]\frametitle{Tópico 02 ---}
\begin{frame}[t]\frametitle{Tópico 02}

  \begin{itemize}
    \justifying{}
    \setlength\itemsep{1em}
    \item Item 1.
    \item Item 2.
    \item Item 3.
  \end{itemize}

\end{frame}



% ------------------------------------------------------------------------------
% ------------------------------------------------------------------------------
% ------------------------------------------------------------------------------
% ------------------------------------------------------------------------------
% ------------------------------------------------------------------------------
\subsection[Exercícios]{Exercícios}\label{subsec:tema04-exercicios}



% \begin{frame}[allowframebreaks,t]\frametitle{Exercícios ---}
\begin{frame}[t]\frametitle{Exercícios}

\begin{itemize}
  \justifying{}
  \setlength\itemsep{1em}
  \item Questão 1.
  \item Questão 2.
  \item Questão 3.
\end{itemize}

\end{frame}


% !TeX encoding = UTF-8
% !TeX spellcheck = pt_BR
% !TeX root = 20231-CODDISC-JOAO-FERREIRA.tex
% ------------------------------------------------------------------------------

% ------------------------------------------------------------------------------
% Template para aulas da UNINASSAU
% Autor: João Ferreira da Silva Júnior
% Versão: 0.4.4 - 2023-07-19
% ------------------------------------------------------------------------------



% ------------------------------------------------------------------------------
% ------------------------------------------------------------------------------
% ------------------------------------------------------------------------------
% ------------------------------------------------------------------------------
% ------------------------------------------------------------------------------
% ------------------------------------------------------------------------------
% ------------------------------------------------------------------------------
% ------------------------------------------------------------------------------
% ------------------------------------------------------------------------------
% ------------------------------------------------------------------------------
\section[Tema 05]{Tema 05}\label{sec:tema05}



{\usebackgroundtemplate{{\transparent{0.1}\includegraphics[width=\paperwidth,height=\paperheight,keepaspectratio]{img/logo/background-branco.png}}}
  % \begin{frame}[allowframebreaks,t]%\frametitle{ ---}
  \begin{frame}[plain]%\frametitle{}

    \vfill

    \centering{}
    \Huge{\textbf{Tema 05}}

    \vfill

  \end{frame}
} % usebackgroundtemplate



% ------------------------------------------------------------------------------
% ------------------------------------------------------------------------------
% ------------------------------------------------------------------------------
% ------------------------------------------------------------------------------
% ------------------------------------------------------------------------------
\subsection[Tópico 01]{Tópico 01}\label{subsec:tema05-topico01}



% \begin{frame}[allowframebreaks,t]\frametitle{Tópico 01 ---}
\begin{frame}[t]\frametitle{Tópico 01}

  \begin{itemize}
    \justifying{}
    \setlength\itemsep{1em}
    \item Item 1.
    \item Item 2.
    \item Item 3.
  \end{itemize}

\end{frame}



% ------------------------------------------------------------------------------
% ------------------------------------------------------------------------------
% ------------------------------------------------------------------------------
% ------------------------------------------------------------------------------
% ------------------------------------------------------------------------------
\subsection[Tópico 02]{Tópico 02}\label{subsec:tema05-topico02}



% \begin{frame}[allowframebreaks,t]\frametitle{Tópico 02 ---}
\begin{frame}[t]\frametitle{Tópico 02}

  \begin{itemize}
    \justifying{}
    \setlength\itemsep{1em}
    \item Item 1.
    \item Item 2.
    \item Item 3.
  \end{itemize}

\end{frame}



% ------------------------------------------------------------------------------
% ------------------------------------------------------------------------------
% ------------------------------------------------------------------------------
% ------------------------------------------------------------------------------
% ------------------------------------------------------------------------------
\subsection[Exercícios]{Exercícios}\label{subsec:tema05-exercicios}



% \begin{frame}[allowframebreaks,t]\frametitle{Exercícios ---}
\begin{frame}[t]\frametitle{Exercícios}

\begin{itemize}
  \justifying{}
  \setlength\itemsep{1em}
  \item Questão 1.
  \item Questão 2.
  \item Questão 3.
\end{itemize}

\end{frame}


% !TeX encoding = UTF-8
% !TeX spellcheck = pt_BR
% !TeX root = 20232-CODDISC-JOAO-FERREIRA.tex
% ------------------------------------------------------------------------------

% ------------------------------------------------------------------------------
% Template para aulas da UNINASSAU
% Autor: João Ferreira da Silva Júnior
% Versão: 0.4.4 - 2023-07-19
% ------------------------------------------------------------------------------



% ------------------------------------------------------------------------------
% ------------------------------------------------------------------------------
% ------------------------------------------------------------------------------
% ------------------------------------------------------------------------------
% ------------------------------------------------------------------------------
% ------------------------------------------------------------------------------
% ------------------------------------------------------------------------------
% ------------------------------------------------------------------------------
% ------------------------------------------------------------------------------
% ------------------------------------------------------------------------------
\section[Listas de Exercícioss]{Listas de Exercícioss}\label{sec:listas}



{\usebackgroundtemplate{{\transparent{0.1}\includegraphics[width=\paperwidth,height=\paperheight,keepaspectratio]{img/logo/background-branco.png}}}
  % \begin{frame}[allowframebreaks,t]%\frametitle{ ---}
  \begin{frame}[plain]%\frametitle{}

    \vfill

    \centering{}
    \Huge{\textbf{Listas de Exercícioss}}

    \vfill

  \end{frame}
} % usebackgroundtemplate



% ------------------------------------------------------------------------------
% ------------------------------------------------------------------------------
% ------------------------------------------------------------------------------
% ------------------------------------------------------------------------------
% ------------------------------------------------------------------------------
\subsection[Lista de Exercícioss 01]{Lista de Exercícioss 01}\label{subsec:listas-lista01}



% \begin{frame}[allowframebreaks,t]\frametitle{Lista de Exercícioss 01 ---}
\begin{frame}[t]\frametitle{Lista de Exercícioss 01}

  \begin{itemize}
    \justifying{}
    \setlength\itemsep{1em}
    \item Item 1.
    \item Item 2.
    \item Item 3.
  \end{itemize}

\end{frame}



% ------------------------------------------------------------------------------
% ------------------------------------------------------------------------------
% ------------------------------------------------------------------------------
% ------------------------------------------------------------------------------
% ------------------------------------------------------------------------------
\subsection[Lista de Exercícioss 02]{Lista de Exercícioss 02}\label{subsec:listas-lista02}



% \begin{frame}[allowframebreaks,t]\frametitle{Lista de Exercícioss 02 ---}
\begin{frame}[t]\frametitle{Lista de Exercícioss 02}

  \begin{itemize}
    \justifying{}
    \setlength\itemsep{1em}
    \item Item 1.
    \item Item 2.
    \item Item 3.
  \end{itemize}

\end{frame}


% !TeX encoding = UTF-8
% !TeX spellcheck = pt_BR
% !TeX root = 20232-CODDISC-JOAO-FERREIRA.tex
% ------------------------------------------------------------------------------

% ------------------------------------------------------------------------------
% Template para aulas da UNINASSAU
% Autor: João Ferreira da Silva Júnior
% Versão: 0.4.4 - 2023-07-19
% ------------------------------------------------------------------------------



% ------------------------------------------------------------------------------
% ------------------------------------------------------------------------------
% ------------------------------------------------------------------------------
% ------------------------------------------------------------------------------
% ------------------------------------------------------------------------------
% ------------------------------------------------------------------------------
% ------------------------------------------------------------------------------
% ------------------------------------------------------------------------------
% ------------------------------------------------------------------------------
% ------------------------------------------------------------------------------
\section[Seminários]{Seminários}\label{sec:seminarios}



{\usebackgroundtemplate{{\transparent{0.1}\includegraphics[width=\paperwidth,height=\paperheight,keepaspectratio]{img/logo/background-branco.png}}}
  % \begin{frame}[allowframebreaks,t]%\frametitle{ ---}
  \begin{frame}[plain]%\frametitle{}

    \vfill

    \centering{}
    \Huge{\textbf{Seminários}}

    \vfill

  \end{frame}
} % usebackgroundtemplate



% ------------------------------------------------------------------------------
% ------------------------------------------------------------------------------
% ------------------------------------------------------------------------------
% ------------------------------------------------------------------------------
% ------------------------------------------------------------------------------
\subsection[Critérios de avaliação]{Critérios de avaliação}\label{subsec:seminarios-criterios}



% \begin{frame}[allowframebreaks,t]\frametitle{Critérios de avaliação ---}
\begin{frame}[t]\frametitle{Critérios de avaliação}

  \begin{itemize}
    \justifying{}
    \setlength\itemsep{1em}
    \item Poderá ser feito individualmente ou em grupo de até 5 (cinco) membros.
    \item Todos do grupo deverão participar tanto da elaboração quanto apresentação, pois todos terão seu desempenho avaliado individualmente e em grupo.
    \item A avaliação do seminário será feita a partir da apresentação para toda a turma.
    \item Serão avaliados os critérios:
    \begin{itemize}
      \justifying{}
      \item Uso do tempo para apresentação de até 20 minutos.
      \item Conteúdo abordado dentro do tema proposto.
      \item Clareza e expressão das ideias.
      \item Domínio do conteúdo apresentado.
      \item Preparo geral da apresentação e processos de pesquisa do tema.
      \item Recursos utilizados, multimídia e criatividade.
    \end{itemize}
  \end{itemize}

\end{frame}



% ------------------------------------------------------------------------------
% ------------------------------------------------------------------------------
% ------------------------------------------------------------------------------
% ------------------------------------------------------------------------------
% ------------------------------------------------------------------------------
\subsection[Temas]{Temas}\label{subsec:seminarios-temas}



% \begin{frame}[allowframebreaks,t]\frametitle{Temas ---}
\begin{frame}[t]\frametitle{Temas}

  \begin{itemize}
    \justifying{}
    \setlength\itemsep{1em}
    \item Item 1.
    \item Item 2.
    \item Item 3.
  \end{itemize}

\end{frame}


% !TeX encoding = UTF-8
% !TeX spellcheck = pt_BR
% !TeX root = 20231-CODDISC-JOAO-FERREIRA.tex
% ------------------------------------------------------------------------------

% ------------------------------------------------------------------------------
% Template para aulas da UNINASSAU
% Autor: João Ferreira da Silva Júnior
% Versão: 0.4.4 - 2023-07-19
% ------------------------------------------------------------------------------



% ------------------------------------------------------------------------------
% ------------------------------------------------------------------------------
% ------------------------------------------------------------------------------
% ------------------------------------------------------------------------------
% ------------------------------------------------------------------------------
% ------------------------------------------------------------------------------
% ------------------------------------------------------------------------------
% ------------------------------------------------------------------------------
% ------------------------------------------------------------------------------
% ------------------------------------------------------------------------------
\section[Projetos]{Projetos}\label{sec:projetos}



{\usebackgroundtemplate{{\transparent{0.1}\includegraphics[width=\paperwidth,height=\paperheight,keepaspectratio]{img/logo/background-branco.png}}}
  % \begin{frame}[allowframebreaks,t]%\frametitle{ ---}
  \begin{frame}[plain]%\frametitle{}

    \vfill

    \centering{}
    \Huge{\textbf{Projetos}}

    \vfill

  \end{frame}
} % usebackgroundtemplate



% ------------------------------------------------------------------------------
% ------------------------------------------------------------------------------
% ------------------------------------------------------------------------------
% ------------------------------------------------------------------------------
% ------------------------------------------------------------------------------
\subsection[Critérios de avaliação]{Critérios de avaliação}\label{subsec:projetos-criterios}



% \begin{frame}[allowframebreaks,t]\frametitle{Critérios de avaliação ---}
  \begin{frame}[t]\frametitle{Critérios de avaliação}

    \begin{itemize}
      \justifying{}
      \setlength\itemsep{1em}
      \item Poderá ser feito individualmente ou em grupo de até 5 (cinco) membros.
      \item Todos do grupo deverão participar tanto da elaboração quanto apresentação, pois todos terão seu desempenho avaliado individualmente e em grupo.
      \item O projeto já concluído e funcionando deverá ser apresentado para toda a turma.
      \item Serão avaliados os critérios:
      \begin{itemize}
        \justifying{}
        \item Uso do tempo para apresentação de até 20 minutos.
        \item Gerenciamento da equipe e consequentes conflitos ou impedimentos.
        \item Processos de desenvolvimento do projeto, modelagem e codificação.
        \item Relato das lições aprendidas e o que mais poderia ser feito.
        \item Clareza e expressão das ideias.
        \item Preparo geral da apresentação.
        \item Recursos utilizados, multimídia e criatividade.
      \end{itemize}
    \end{itemize}

  \end{frame}



% ------------------------------------------------------------------------------
% ------------------------------------------------------------------------------
% ------------------------------------------------------------------------------
% ------------------------------------------------------------------------------
% ------------------------------------------------------------------------------
\subsection[Projeto 01]{Projeto 01}\label{subsec:projetos-projeto01}



% \begin{frame}[allowframebreaks,t]\frametitle{Projeto 01 ---}
\begin{frame}[t]\frametitle{Projeto 01}

  \begin{itemize}
    \justifying{}
    \setlength\itemsep{1em}
    \item Item 1.
    \item Item 2.
    \item Item 3.
  \end{itemize}

\end{frame}



% ------------------------------------------------------------------------------
% ------------------------------------------------------------------------------
% ------------------------------------------------------------------------------
% ------------------------------------------------------------------------------
% ------------------------------------------------------------------------------
\subsection[Projeto 02]{Projeto 02}\label{subsec:projetos-projeto02}



% \begin{frame}[allowframebreaks,t]\frametitle{Projeto 02 ---}
\begin{frame}[t]\frametitle{Projeto 02}

  \begin{itemize}
    \justifying{}
    \setlength\itemsep{1em}
    \item Item 1.
    \item Item 2.
    \item Item 3.
  \end{itemize}

\end{frame}




% ------------------------------------------------------------------------------
% ------------------------------------------------------------------------------
% ------------------------------------------------------------------------------
% ------------------------------------------------------------------------------
% ------------------------------------------------------------------------------
% ------------------------------------------------------------------------------
% ------------------------------------------------------------------------------
% ------------------------------------------------------------------------------
% ------------------------------------------------------------------------------
% ------------------------------------------------------------------------------
\section[Agradecimento]{Agradecimento}\label{sec:agradecimentos}



{\usebackgroundtemplate{{\transparent{0.2}\includegraphics[width=\paperwidth,height=\paperheight,keepaspectratio]{img/logo/background-amarelo.png}}}
  \begin{frame}[plain,t]

    \vfill
    \vspace*{0.15\paperheight}

    \begin{center}
      \Huge{\textbf{Obrigado!}}
    \end{center}

    \vspace*{0.15\paperheight}
    \vfill

    \begin{block}{\professor}
      \justifying{}
      \begin{itemize}
        \justifying{}
        \setlength\itemsep{1em}
        \item \href{mailto:\professoremail}{\professoremail}
        \item \href{\professorlattes}{\professorlattes}
        \item \href{professorsite}{\professorsite}
      \end{itemize}
    \end{block}

  \end{frame}
} % usebackgroundtemplate



% ------------------------------------------------------------------------------
% ------------------------------------------------------------------------------
% ------------------------------------------------------------------------------
% ------------------------------------------------------------------------------
% ------------------------------------------------------------------------------
% ------------------------------------------------------------------------------
% ------------------------------------------------------------------------------
% ------------------------------------------------------------------------------
% ------------------------------------------------------------------------------
% ------------------------------------------------------------------------------
\section[Referências]{Referências}\label{sec:referencias}



{\usebackgroundtemplate{{\transparent{0.1}\includegraphics[width=\paperwidth,height=\paperheight,keepaspectratio]{img/logo/background-amarelo.png}}}
  \begin{frame}[plain,allowframebreaks,t]%\frametitle{Referências ---}
  % \begin{frame}[plain,t]\frametitle{Referências}

    \begin{center}
      \Huge{\textbf{Referências}}
    \end{center}

    % \bibliographystyle{apalike}
    \bibliography{library}

  \end{frame}
} % usebackgroundtemplate



% ------------------------------------------------------------------------------
% ------------------------------------------------------------------------------
% ------------------------------------------------------------------------------
% ------------------------------------------------------------------------------
% ------------------------------------------------------------------------------
% ------------------------------------------------------------------------------
% ------------------------------------------------------------------------------
% ------------------------------------------------------------------------------
% ------------------------------------------------------------------------------
% ------------------------------------------------------------------------------
\section*{Capa Final}\label{sec:capa-final}



{\usebackgroundtemplate{\includegraphics[width=\paperwidth,height=\paperheight,keepaspectratio]{img/logo/background-branco.png}}
  \begin{frame}[plain]

    ~

  \end{frame}
} % usebackgroundtemplate



\end{document}
